\weeknum{2}
\topic[Physical Properties]{Physical Properties, Mixing Laws, and Asking the Right Questions}

\workshopstart

\vspace{-0.75em}
\noindent\textbf{Duration:} 3 hours \hfill \textbf{Setting:} Workshop room

\section*{Preparation}

\begin{description}
  \item[Pre-reading:] Course Notes Ch. 2
  \item[Lecture videos:]
  \item[Notes:]
\end{description}

\section*{Purpose}

This workshop develops intuition for how physical properties (density $\rho$, thermal conductivity $k$, thermal diffusivity $\kappa$, bulk modulus $K_T$, shear modulus $\mu$, seismic velocities $V_P, V_S$, magnetic susceptibility $\chi$, and heat production $A$) depend on state variables (pressure $P$, temperature $T$, composition, porosity/cracks), and how geometry and fabric determine appropriate mixing relations (arithmetic, geometric, harmonic; Voigt--Reuss--Hill). Students practice cross-disciplinary collaboration and formulate the key questions a practicing geophysicist must ask a geologist before modeling or surveying.

\section*{Learning Outcomes}
\begin{itemize}
  \item Select and justify mixing laws based on rock fabric (random granular, layered/foliated, fractured) and target property (bulk vs. transport vs. elastic).
  \item Reason about how $P$--$T$--composition/porosity shift properties and create anisotropy.
  \item Identify which properties most strongly influence gravity, magnetic, and heat flow interpretations.
  \item Elicit geological information via well-posed questions to constrain non-uniqueness.
\end{itemize}

\section*{Session Flow (\timelinelength\ min)}

Create mixed groups (4--5 students each) with both geology and physics/engineering backgrounds.

\timeline[slot]{Warm-up --- What does geophysics see?}{10}

Much like our senses, different fields are sensitive to different things.  Are geophysical fields sensitive to individual minerals or bulk rock propoerties?  Are they sensitive to different scales?

\timeline[slot]{Mini-lecture and Q\&A}{20}
\begin{itemize}
  \item Bulk-additive vs. pathway-limited vs. elastic-interaction properties.
  \item Mixing relations: arithmetic (parallel), harmonic (perpendicular), geometric (random granular), Voigt--Reuss--Hill (elastic bounds).
  \item Scale-dependent anisotropy (micro/meso/macro fabrics).
\end{itemize}

\timeline[slot]{Live Demonstration --- ``Crowd Conductivity''}{15}

\textbf{Aim:} Intuitively illustrate why different mixing laws arise and how anisotropy appears.
\begin{enumerate}
  \item \textbf{Random scattering} (students distributed irregularly): volunteer transits across room. Multiple semi-obstructed pathways $\Rightarrow$ geometric-like effective behavior.
  \item \textbf{Parallel corridors} (lanes aligned with motion): fastest transit $\Rightarrow$ arithmetic mean (parallel pathways).
  \item \textbf{Perpendicular barriers} (rows perpendicular to motion): slowest transit $\Rightarrow$ harmonic mean (series resistances).
  \item \textbf{Anisotropy}: compare transit parallel vs. perpendicular to fabric.
\end{enumerate}
\textit{Optional: time a handful of transits and note distribution shape differences.}

\timeline[slot]{Simulation Walkthrough}{15}

Builds on crowd conductivity with computation.  We'll examine a pre-built Python visualization of transit times through three media (random, parallel, perpendicular) and the resulting histogram shapes. Relate to effective properties and mixing choices.

\timeline[item]{Rock Fabric and Mixing Laws}{20}

Distribute a set of photos (thin section, hand sample/outcrop close-up, outcrop/aerial). For each image, groups identify: geometry (random granular/layered/foliated/fractured), appropriate mixing law for density, and conductivity, whether the material is isotropic at that scale, and how conclusions change with scale.

\timeline[item]{When to Estimate Physical Properties}{45}

Each group tackles two scenarios (\~20 min each: 12 min discussion, 5 min notes, 3 min question drafting). Use the dedicated scenario handouts. Outcomes must include: (i) property/mixing-law decisions and (ii) 2--3 \emph{questions to ask the geologist} before committing to a model/survey.

\timeline[demo]{Path Simulation}{30}

In this computer demonstration, we'll explore how the arrangement of properties affects the path and rate of travel across a medium.  While not a true simulator of physical properties, the results will show how various geometries of multicomponent systems affect the path and travel-time.

\timeline[slot]{Wrap-up}{10}

\emph{What does geophysics see?} Geophysics responds to bulk physical properties, not individual minerals. The appropriate mixing law depends on the rock fabric and the property being estimated. Can you name a physical property for each of the geophysical methods introduced so far (gravity, magnetics, heat flow)?

\timeline[slot]{Preview: Inversion}{10}

In geophysics, we typically measure fields and infer the underlying physical properties. In this course, however, we mostly focus on the \emph{forward problem}: computing fields from a known distribution of properties.

In practice, we usually face the opposite challenge: the \emph{inverse problem}.
Inverse problems are difficult for several reasons:
\begin{itemize}
  \item they are often \emph{underdetermined}, with fewer observations than unknowns;
  \item \emph{non-unique}, meaning different models can produce the same data; and
  \item \emph{unstable}, where small changes in the model can produce large changes in the predicted field.
\end{itemize}
In addition, many geophysical relationships are non-linear, which adds further complexity.

Next week, we'll begin looking at how inversion methods address these challenges.

\workshopend
\notepage
\notepage

\subimport{activities/}{activity_photo_mixing}
\subimport{activities/}{activity_scenarios}
\subimport{activities/}{demo_path_simulation}